\section{Potential Sources of Data Loss in the HTA Signal Chain}

The HTA module converts mechanical load into an electrical command that ultimately drives a motor. Along this path—from strain transfer in the structure, through bridge sensing and digitization, to digital processing and re-generation of an analogue command—information can be attenuated, distorted, delayed, or dropped. The following hypothesis outlines where and why this can occur.

\textbf{Mechanical–to–electrical transduction.} Data quality begins at the strain gauges. Imperfect bonding, local bending not representative of the global load, gauge misalignment, or partial debonding causes under-transfer of strain, i.e., real load changes do not fully appear at the gauge. Leadwire movement and cable microphonics can add spurious strain. Creep and hysteresis in the adhesive stack, as well as temperature gradients between gauges, introduce slow drift that obscures small signals.

\textbf{Wheatstone bridge and analogue wiring.} Series lead resistance and unequal gauge self-heating unbalance the bridge and compress dynamic range around an offset. Insufficient shielding or twisted-pair routing lets 50/60 Hz fields and motor EMI couple into the sense lines, masking high-frequency details. If the bridge excitation sags with battery voltage, the scale factor varies over time (gain error), which looks like loss or gain of data depending on state of charge.

\textbf{HX711 acquisition.} Although the HX711 offers high resolution, information can be lost if the input saturates (over-range at high gain), if conversions are read before the device is ready, or if the selected sample rate under-samples the true bandwidth of the load signal, causing aliasing. Conversion jitter and settling after gain/channel changes momentarily corrupt the stream. Poor power decoupling at the HX711 or a noisy ground shared with the motor return injects quantization steps that are not related to the load.

\textbf{Power integrity and grounding.} The LM358 stages and the bridge are sensitive to supply noise. H-bridge commutation, motor current steps, or converter ripple can bounce the common ground; any movement of the analogue reference effectively subtracts information from the signal. Inadequate bypass capacitors, long return loops, or mixing high-current and low-level returns lead to repeatable yet hard-to-see data erosion.

\textbf{Digital filtering and control.} A moving-average filter improves SNR but also attenuates rapid, legitimate changes and introduces delay; fast transients can be “averaged away.” The derivative term in the PD branch amplifies residual noise and, when clipped or rate-limited, discards peaks. Finite numeric precision (e.g., integer rounding, overflow) and sampling-period jitter smear small variations into the baseline.

\textbf{Command re-generation (DAC path).} The MCP4725’s 12-bit quantization limits the smallest representable step. Code rounding and I\textsuperscript{2}C retries (or missed ACKs on a marginal bus) momentarily hold the previous code—perceived as lost updates. Downstream, the LM358 has finite output swing and slew rate; if asked to jump faster or higher than its limits, it flattens the waveform, removing amplitude and edge information. Any intentional RC smoothing further trades detail for noise reduction.

\textbf{Actuator coupling and back-injection.} The H-bridge and motor can feed disturbances back into the analogue ground and supplies. Without adequate flyback diodes, snubbing, and layout separation, these spikes manifest as dropped or distorted ADC samples during switching windows.

\textbf{Encoder feedback (validation path).} Missed or extra counts arise from marginal slot penetration, contamination on the disc, or inadequate interrupt bandwidth at higher speeds. Debounce or digital filtering can also suppress legitimate pulses. When the feedback under-reports speed, comparisons against the commanded signal misleadingly suggest loss on the input side.

\textbf{Calibration and drift.} Offset and span derived from the two-point procedure can be biased if the “standing” or “hanging” conditions are not stable. Temperature changes between calibration and operation shift both bridge and amplifier gains; untracked drift behaves like slow data leakage.

\textbf{Presentation and logging.} The phone dashboard and oscilloscope views are typically down-sampled for rendering. If the display/update rate is lower than the acquisition rate, intermediate samples are dropped by design; the visualization then shows an apparently “thinner” signal even when the raw stream is intact.

\textit{Implication.} The dominant risks concentrate at (i) strain transfer and bridge integrity, (ii) power/ground quality shared with the drive, and (iii) rate/quantization choices in the HX711–filter–DAC chain. Mitigations follow directly: ensure robust bonding and temperature pairing of gauges; isolate analogue and power returns with star-grounding and generous local decoupling; choose sample rates above twice the signal bandwidth; document filter delay; validate DAC/I\textsuperscript{2}C error handling; and verify encoder throughput as an independent check on the commanded dynamics.
